\phantomsection
\section*{ВВЕДЕНИЕ}
\addcontentsline{toc}{section}{Введение}

Подшипники скольжения относятся к числу базовых элементов нефтегазового оборудования: они применяются в магистральных и подпорных насосных агрегатах, в центробежных и винтовых компрессорах, в редукторах буровых установок и приводах турбомашин. Условия их работы~--- высокие удельные нагрузки на цапфу, большие окружные скорости и непрерывный режим эксплуатации~--- делают вопрос ресурса и КПД таких узлов одним из определяющих для надёжности всего технологического оборудования.

Гидродинамический режим работы подшипника достигается, когда в клиновом зазоре между валом и вкладышем за счёт вязкого увлечения масла формируется поле давления, способное полностью разделить трущиеся поверхности. На переходных режимах~--- при пуске, останове и кратковременных перегрузках~--- этот режим может срываться, и поверхности контактируют через тонкий граничный слой, что приводит к ускоренному износу. Любые конструктивные или технологические решения, повышающие устойчивость гидродинамического режима и расширяющие область его существования, представляют непосредственный практический интерес.

Одно из таких решений~--- регулярное текстурирование рабочей поверхности вкладыша, то есть формирование на ней массива микроуглублений известной формы и плотности расположения. Технологически наиболее воспроизводимым способом нанесения текстур является лазерная обработка: импульсное излучение испаряет материал с микронной точностью по глубине, а контроль формы пятна позволяет получать углубления различных профилей~--- эллипсоидальных, цилиндрических, конических, сферических, плоскодонных и комбинированных. Многочисленные экспериментальные и расчётные исследования последних двух десятилетий показывают, что грамотный выбор параметров текстуры приводит к снижению коэффициента трения на 10-30\% и одновременному росту нагрузочной способности.

Физическая основа эффекта~--- локальный микрогидродинамический клиновой зазор, формирующийся на входной по ходу вращения кромке каждого углубления. На этой кромке возникает дополнительный подъём давления, на выходной кромке~--- падение, ограниченное снизу нулём из-за кавитации. В сумме каждое углубление вносит положительный вклад в подъёмную силу. Дополнительно углубления выполняют функции локальных резервуаров смазки на пусковых режимах и ловушек для частиц износа.

Количественное исследование подобных эффектов проводится на основе уравнения Рейнольдса~--- эллиптического дифференциального уравнения для давления в тонком вязком слое. В работе уравнение решается численно методом конечных разностей; итерации проводятся по схеме Гаусса--Зейделя с верхней релаксацией, позволяющей ускорить сходимость на крупных сетках. Результат расчёта~--- двумерное поле давления, по которому интегрированием получают нагрузочную способность, коэффициент трения и расход смазки.

Практическая ценность темы определяется тем, что текстурирование вкладышей легко встраивается в существующий технологический маршрут изготовления подшипников и не требует пересмотра сопрягаемых деталей или конструкции узла в целом. Это делает данное решение применимым в задачах модернизации действующего оборудования.

Цель работы~--- получение и сопоставительный анализ трибологических характеристик радиального подшипника скольжения с дискретным микрорельефом методом численного моделирования.

В рамках поставленной цели решаются следующие задачи:
\begin{enumerate}
    \item изучить теоретические основы гидродинамической смазки и обобщить литературные данные о работе подшипников скольжения с регулярным микрорельефом;
    \item построить численную схему решения уравнения Рейнольдса методом конечных разностей с итерациями по Гауссу--Зейделю;
    \item рассчитать поля давления и интегральные характеристики гладкого и текстурированного подшипников во всём диапазоне рабочих эксцентриситетов;
    \item провести сопоставительный анализ полученных результатов и количественно оценить эффект применения микрорельефа.
\end{enumerate}

Объектом исследования является радиальный подшипник скольжения с \varDimpleTypeGenitive{} микрорельефом рабочей поверхности (вариант~\varNumber, геометрия~\varGeomLetter).

Предметом исследования служат закономерности изменения нагрузочной способности, коэффициента трения и объёмного расхода смазки при нанесении конических углублений на поверхность вкладыша.

Методы исследования: численное решение уравнения Рейнольдса методом конечных разностей, итерационная процедура Гаусса--Зейделя с верхней релаксацией, сопоставительный анализ интегральных трибологических показателей.

Структурно работа состоит из введения, трёх глав, заключения, списка использованных источников и приложения. Первая глава содержит теоретический обзор гидродинамической смазки и анализ литературы по лазерному текстурированию поверхностей трения. Во второй главе формулируется математическая модель и описывается применённая численная схема. В третьей главе приводятся расчётные результаты и проводится их сопоставительный анализ.
